\section{Institutional Background}
\label{sec:institutional}

This section describes the institutional framework governing health litigation and pharmaceutical procurement in Brazil, with a focus on the state of S\~ao Paulo.

\subsection{Health as a Constitutional Right and the Rise of Litigation}
\label{sec:health_right}

The 1988 Brazilian Constitution created the Unified Health System (SUS), a universal public health system jointly funded and administered by federal, state, and municipal governments. The SUS maintains a formulary of approved medications and procedures, periodically updated to reflect medical advances subject to cost-effectiveness criteria. This formulary represents the society's chosen balance between universal coverage and fiscal sustainability \citep{castro2019, soares2019}.

However, Brazilian courts have adopted a strict literal interpretation of Article~196 of the Constitution, which states that ``health is a right of all and a duty of the state.'' Under this interpretation, any individual can obtain virtually any medication through litigation, irrespective of cost or inclusion on the SUS formulary. The range extends from acetylsalicylic acid to galsulfase (for mucopolysaccharidosis type~VI, with annual treatment costs of approximately US\$400,000).

Several factors make health litigation particularly prevalent. Access to the legal system is inexpensive---requiring only a prescription and a public defender---and success rates are exceptionally high: approximately 85\% of first-instance claims in S\~ao Paulo are granted \citep{cnj2019judicializacao}, with even higher rates on appeal. These conditions create strong incentives for individuals to seek medicines through the courts rather than through standard SUS channels.

The growth in health litigation has been dramatic. Nationally, first-instance health court orders reached approximately 96,000 in 2017, representing a 130\% increase over 2008 \citep{cnj2019judicializacao}. In S\~ao Paulo, the increase was 913\%, from 2,317 to 23,465 annual lawsuits between 2008 and 2017. Court orders span a wide geographic distribution across the state's municipalities (Figure~\ref{fig:distribution}) and encompass approximately 2,760 distinct items ordered at least once per year.

\begin{figure}[ht]
  \centering
  \includegraphics[width=0.8\textwidth]{figure1.png}
  \caption{Distribution of Health Litigation Cases Across Municipalities in S\~ao Paulo (2008--2017)}
  \label{fig:distribution}
  \par\smallskip\footnotesize\textit{Source:} S-CODES (SES/SP) and CNJ/INSPER (2019).
\end{figure}

The fiscal impact is substantial. S\~ao Paulo's annual public health budget was approximately US\$5.9~billion in 2018 (about 10\% of the state's total budget). Public spending on litigated health items alone amounted to nearly 5\% of this budget, or approximately US\$300~million annually.

\subsection{Public Procurement under Judicial Pressure}
\label{sec:procurement_pressure}

The S\~ao Paulo State Department of Health (SES/SP) manages health policy implementation through 99 decentralized public buyer units (PBUs) distributed across the state (Figure~\ref{fig:pbu}).

\begin{figure}[ht]
  \centering
  \includegraphics[width=0.8\textwidth]{figure2.png}
  \caption{Public Buyer Units (SES/SP)}
  \label{fig:pbu}
\end{figure}

Federal Law 8,666/1993 establishes the general framework for public procurement, requiring that all government purchases follow a formal tendering process organized in three sequential phases:
\begin{enumerate}
  \item \textbf{Internal phase (planning):} PBUs identify needs, specify items, set quantities and reference prices, choose the auction format, and publish a public notice.
  \item \textbf{External phase (bidding):} Firms compete through a chosen competitive procedure---either sealed-bid (\textit{convite}, up to R\$176,000) or multi-round descending auction (\textit{preg\~ao}, no value limit).
  \item \textbf{Delivery phase:} The winning firm delivers the procured items.
\end{enumerate}

For ordinary purchases, the internal phase typically spans 30 to 180~days, allowing PBUs to aggregate demand, research market prices, and set competitive reference prices. Court orders fundamentally disrupt this process. Almost all court decisions (99.94\%) are enforced as preliminary injunctions with delivery deadlines of 1 to 10~days, compressing the internal phase to roughly one-third of the ordinary timeline.

This compression affects procurement in several ways. First, PBUs lose the ability to aggregate demand across units and over time, reducing their bargaining leverage from bulk purchasing. Second, reference prices must be set quickly with limited market research, often resulting in higher reservation prices. Third, approximately 75\% of litigated items are not on the SUS formulary, meaning PBUs have less experience procuring them. Fourth, roughly 65\% of court orders involve ``packages'' combining SUS and non-SUS items, further complicating procurement planning.

\subsection{The Sanction Channel: Administrative vs.\ Litigated Purchases}
\label{sec:sanction_channel}

A critical institutional feature for our identification strategy is the distinction between \textit{litigated} and \textit{administrative} urgent purchases. Both types face identical planning constraints---compressed timelines, small quantities, and diverted budget resources. However, they differ sharply in the consequences of procurement failure.

For litigated purchases, failure to comply with a court order exposes public officials to severe sanctions: substantial fines (sometimes disproportionate to the purchase value), administrative and criminal liability, and seizure or blocking of public funds. These penalties are public information, as the court order must be referenced in the tender notice.

Since 2009, the SES/SP has attempted to mitigate these costs through administrative requests, a mechanism allowing the government to negotiate supply directly with individuals before litigation occurs. Administrative requests undergo evaluation by a scientific committee using evidence-based medicine criteria. Crucially, failure to complete an administrative purchase carries no penalties for public officials. Between 2009 and 2019, administrative requests generated approximately 9,700 purchases, or about 6.5\% of the total from court orders.

This institutional distinction provides a natural comparison group for isolating the ``under the gun'' effect---the additional procurement cost attributable to the threat of judicial sanctions---since administrative and litigated purchases share all planning and execution constraints but differ only in the penalty regime.

\begin{figure}[ht]
  \centering
  \includegraphics[width=0.8\textwidth]{figure3.pdf}
  \caption{Types of Purchases}
  \label{fig:purchase_types}
\end{figure}

\subsection{The Bidding Process}
\label{sec:bidding}

Both ordinary and urgent purchases are conducted through the \textit{Bolsa Eletr\^onica de Compras} (BEC), S\~ao Paulo state's electronic procurement platform, mandatory for all common goods and services since 2007. The BEC handles substantial volumes: approximately US\$1.7~billion in SES/SP transactions in 2018, and over R\$7.4~billion since 2009.

Two competitive procedures are used. In sealed-bid tendering (\textit{convite}), firms submit proposals by a deadline; the lowest qualifying bid wins. In multi-round descending auctions (\textit{preg\~ao}), an initial sealed-bid phase is followed by a 20-minute reverse auction where firms observe the current lowest bid and can submit lower bids. If a bid is placed in the final four minutes, the auction extends by four additional minutes. After the auction, a negotiation phase allows further price reductions.

The key implication for our analysis is that all purchase types use the same electronic platform and competitive procedures. The variation we exploit arises exclusively from differences in planning conditions and penalty regimes, not from differences in the auction mechanism itself.
